
We designed an experimental pipeline to test whether geometric properties of hidden state representations correlate with semantic entropy in large language models. The research question was whether spectral features (participation ratio, eigenvalue decay, condition number) extracted from layers 24-31 during single forward passes achieve Spearman correlation magnitude exceeding 0.4 with semantic entropy on TruthfulQA factual questions.

The implementation achieved completion (approximately 1,200 lines of code across five modules) and validated preliminary steps (data loading, model initialization). However, hidden state extraction for 7B parameter models proved computationally infeasible on LIGHT tier resources, consuming more than 10 hours without completion for 246 test examples. This represents a 20-fold underestimate of planning complexity scores.

The hypothesis remains untested. No correlation measurements were obtained. The research question is unanswered. We quantify the computational barrier: 246 examples × 8 layers × 4096 dimensions × 2 bytes ≈ 15.4 MB tensor operations with extrapolated per-example cost of approximately 12 minutes.

This finding indicates that geometric uncertainty quantification for 7B-scale models requires either: (1) optimized extraction using FlashAttention-2, vLLM, or similar inference libraries, (2) streaming/checkpointing strategies to reduce memory overhead, (3) resource scaling to higher computational tiers, or (4) model downsizing to sub-1B parameters where extraction becomes tractable on standard resources.

The failure documents a computational constraint others pursuing geometric approaches will encounter and provides quantitative data for resource allocation decisions. The methodology is sound, the implementation is complete, and the computational requirements are now specified. The hypothesis awaits validation with tractable model scales, optimized extraction infrastructure, or upgraded computational resources.
